\documentclass[10pt,a4paper]{amsart}
\usepackage[margin=19mm]{geometry}
\usepackage{amsmath,amssymb,mathtools}
\usepackage[scr=rsfs,cal=euler]{mathalfa}
\usepackage{xfrac}
\usepackage[unicode,hidelinks,bookmarks=false]{hyperref}
\newcommand{\ie}{i.e.\ }
\newcommand{\cf}{cf.\ }
\newcommand{\resp}{respectively\ }
\newcommand{\Subsection}{\subsection}
\providecommand{\nobreakdash}{\nobreak\hbox{-}}
\setlength{\emergencystretch}{2em}
\multlinegap0pt
\usepackage{mathtools}
\usepackage{amssymb}
\usepackage{dsfont}
\usepackage{tikz}
\usepackage[shortlabels]{enumitem}
\newtheorem{theorem}{Theorem}
\newtheorem{proposition}{Proposition}
\DeclareMathOperator{\Ext}{Ext}
\newcommand{\F}{\mathbb F}
\DeclareMathOperator{\Hm}{Hom}
\newcommand{\Q}{\mathbb Q}
\newcommand{\Z}{\mathbb Z}
\renewcommand{\ge}{\geqslant}
\def\k{\mathbf{k}}
\title{On formality of diagrams of Eilenberg-MacLane spaces}
\author{Grigory Solomadin, Antoine Touz\'e}
\date{Source: arXiv:2604.22435v1 (CC BY 4.0)}
\begin{document}
\maketitle

\noindent\emph{Source and licence.} On formality of diagrams of Eilenberg-MacLane spaces. Version arXiv:2604.22435v1 (CC BY 4.0), distributed by arXiv under the Creative Commons Attribution 4.0 International licence, http://creativecommons.org/licenses/by/4.0/, as recorded in the arXiv licence metadata for this exact version and shown on the abstract page https://arxiv.org/abs/2604.22435v1. Full licence notice: \texttt{licenses/arxiv-2604.22435-CC-BY-4.0.txt}.\par\smallskip

\noindent\textit{Editorial scope.} Counted unit: the article's theorem thm:nonformality (source0.tex lines 555--557). The article's complete original proof of it is delivered with it (source0.tex lines 622--649, one \texttt{proof} environment of 28 lines, which the article places away from the statement). No mathematics is written, repaired or re-derived: the statement and the proof are the article's own lines, in the article's own notation, and no retained line is substituted or re-flowed.  Retained uncounted as the source-local context the proof needs: lines 601--610; lines 616--618.  Nothing was omitted from any retained span, so the package carries no omission record.  No retained line contains a citation, so the wrapper carries no bibliography and no bibliography entry.  No retained line contains a figure environment, an image-inclusion command or drawing code, so no image is delivered and the package declares no asset dependency.  Editorial formatting only, in addition: an AMS \texttt{amsart} preamble that restates the article's own declarations and macros used by the retained lines, namely the environments \texttt{theorem}, \texttt{proposition} on separate counters, together with the standard packages those lines need.  The retained environments print in this document as Theorem 1, Proposition 1 in order of appearance, whereas the article prints the same items with the numbering of its own full document.  The complete native source of the article as distributed in the arXiv e-print for this exact version is bundled unchanged as \texttt{sources/199092/source0.tex}.
\begin{proposition}\label{prop:Extcomputation1}
	Let $i$ and $k$ be two nonnegative integers. We have 
	\[
		\Ext^i_{\mathcal{F}(\langle C\rangle,\k)}(\overline{P}^{\otimes k},A)=
		\begin{cases}
			\F & \text{if $k=1$ and $i=0$,}\\
			0  & \text{otherwise.}
		\end{cases}
	\]
\end{proposition}

\begin{proposition}\label{prop:Extcomputation2}
	For all positive $n$, the $\F$-vector space $\bigoplus_{i,j\ge 0}\Ext^i_{\mathcal{F}(\langle C\rangle,\k)}(H_j(n,\k),A)$ has dimension greater than one.
\end{proposition}

\begin{theorem}\label{thm:nonformality}
	Assume that $\k$ is not a $\Q$-algebra, and let $C$ be a cyclic group such that $C\otimes_{\Z}\k\ne 0$. Then for all positive $n$ the functor $\overline{B}^n\k[-]:\langle C\rangle\to \mathrm{DGMod}_\k$ is not formal.
\end{theorem}

\begin{proof}[Proof of theorem~\ref{thm:nonformality}]
	We proceed by contradiction, so we assume that $\overline{B}^n\k[-]$ is formal. If we regard $\overline{B}^n\k[-]$ as a complex in $\mathcal{F}(\langle C\rangle,\k)$, then formality means that it is connected to the complex $H(n,\k)=\bigoplus_{i\ge 0}H_i(n,\k)$ (with zero differential) by a zig-zag of quasi-isomorphisms of complexes in $\mathcal{F}(\langle C\rangle,\k)$. Thus, if $J^A$ denotes an injective resolution of $A$ in $\mathcal{F}(\langle C\rangle,\k)$, then the totalization of the bicomplex of $\k$-modules
	\[
		T:=\Hm_{\mathcal{F}(\langle C\rangle,\k)}(\overline{B}^n\k[-],J^A)
	\]
	is connected by a zig-zag of quasi-isomorphisms in $\mathrm{DGMod}_\k$ to the total complex of $\k$-modules
	\[
		T':=\Hm_{\mathcal{F}(\langle C\rangle,\k)}(H(n,\k),J^A)\;.
	\]
	
	Let us compute the homology of $T$. Let $\overline{B}^n\k[-]_s$ denote the object of degree $s$ of the complex  $\overline{B}^n\k[-]$. There is a first quadrant spectral sequence:
	\[
		E_1^{st}=\Ext^t_{\mathcal{F}(\langle C\rangle,\k)}(\overline{B}^n\k[-]_s,A)\Rightarrow H^{s+t}(T)\;.
	\]
	As a graded object, the reduced normalized bar complex of an augmented dg-algebra $A$ is isomorphic to $\bigoplus_{k\ge 0}\overline{A}[1]^{\otimes k}$ with the convention that $\overline{A}[1]^{\otimes 0}:=\k$, 
	thus we see by induction on $n$ that:
	\[ 
		\overline{B}^n\k[-]_s = 
		\begin{cases}
			\k & \text{if $s=0$,}\\
			0 & \text{if $0<s<n$,}\\
			\overline{P} & \text{if $s=n$,}
		\end{cases}
	\]
	and for $s>n$, $\overline{B}^n\k[-]_s$ is a direct sum of copies of functors $\overline{P}^{\otimes k}$ with $k>1$. Therefore by proposition~\ref{prop:Extcomputation1}, we have $E_1^{st}\simeq\F$ if $(s,t)=(n,0)$ and zero otherwise. As a consequence, the spectral sequence collapses and the total homology of $T$ is isomorphic to $\F$.
	
	Now the total homology of $T'$ is $\bigoplus_{i,j\ge 0}\Ext^i_{\mathcal{F}(\langle C\rangle,\k)}(H_j(n,\k),A)$, which is an $\F$-vector space of dimension greater than one by proposition~\ref{prop:Extcomputation2}. This contradicts the fact that $T$ and $T'$ have isomorphic homology. Hence $\overline{B}^n\k[-]$ cannot be formal.
\end{proof}

\end{document}
